% Teil II: der durchgerechnete Weg.

% Gedankenschritt: Summanden erkennen, Startwert anpassen; Teleskopsumme
% nutzen, um große Summen ohne Gliedweise-Berechnung auszuwerten.

\subsection*{Gegeben und gesucht}
(a) Summe mit $\Sigma$ schreiben; (b) Wert für $n=1,2,3$; (c) Summe
ausschreiben und Wert berechnen.

\subsection*{a) $2\cdot3+3\cdot4+\dots+n(n+1)$}
Der $k$-te Summand ist $k(k+1)$. Der erste Term entspricht $k=2$
(denn $2\cdot3$), der letzte $k=n$:
\begin{equation}
  2\cdot3+3\cdot4+\dots+n(n+1)=\sum_{k=2}^{n}k(k+1).
\end{equation}
\textbf{Gedankenschritt:} Start bei $k=2$, weil $k=1$ den Term $1\cdot2$
liefern würde, der nicht vorkommt.

\subsection*{b) $\sum_{k=0}^{n}\frac{(3k-1)^2}{(k+1)!}$}
Mit $(k+1)!=(k+1)k!$ gliedweise auswerten (für kleine $n$; der
entscheidende Trick ist die Fakultätsbeziehung, damit man Brüche
sauber ergänzt):
\begin{align}
  n=1:&\quad \frac{(-1)^2}{1!}+\frac{2^2}{2!}=1+2=3, \\
  n=2:&\quad 3+\frac{5^2}{3!}=3+\frac{25}{6}=\frac{43}{6}, \\
  n=3:&\quad \frac{43}{6}+\frac{8^2}{4!}=\frac{43}{6}+\frac{64}{24}
        =\frac{43}{6}+\frac{8}{3}=\frac{59}{6}.
\end{align}

\subsection*{c) $\sum_{i=1}^{1000}(i^2-(i-1)^2)$}
Ausschreiben:
\begin{equation}
  (1^2-0^2)+(2^2-1^2)+(3^2-2^2)+\dots+(1000^2-999^2).
\end{equation}
\textbf{Der entscheidende Gedankenschritt: Teleskopieren.} Jeder negative
Term $-i^2$ erscheint als $-(i-1)^2$ im nächsten Glied und hebt sich mit
dem positiven $i^2$ des Folgeglieds weg. Übrig bleiben nur der erste und
der letzte "$a$-Wert``:
\begin{equation}
  =1000^2-0^2=1000000 .
\end{equation}

\begin{ergebnis}
\begin{equation*}
  \text{a)}\ \sum_{k=2}^{n}k(k+1),\qquad
  \text{b)}\ n=1{:}\ 3,\ n=2{:}\ \tfrac{43}{6},\ n=3{:}\ \tfrac{59}{6},\qquad
  \text{c)}\ 1000000.
\end{equation*}
\end{ergebnis}

\begin{kontrolle}
\textbf{b) nachrechnen über die Summanden:} $n=2$: $(3k-1)^2/(k+1)!$:
$k=0$: $1$, $k=1$: $2$, $k=2$: $25/6$, Summe $\frac{43}{6}$. \checkmark

\textbf{c) Grenzfall / Generalisierung:} Allgemein hat
$\sum_{i=1}^{N}(i^2-(i-1)^2)=N^2$ (Teleskop). Der Wert ließe sich auch
über $\sum(2i-1)=N^2$ bekommen (die Tonglieder $i^2-(i-1)^2=2i-1$).
Für $N=1000$ ist $N^2=10^6$, plausibel.

\textbf{Stichprobe:} Für $N=4$: $1+3+5+7=16=4^2$; die Teleskopsumme
liefert $4^1-0=16$. Passt.
\end{kontrolle}

\begin{merksatz}
Teleskopsummen kürzen sich fast vollständig --- es bleiben nur der erste
und der letzte Term; beim $\Sigma$-Schreiben den passenden Startindex
wählen. Fakultätsummen nutzt man über $(k+1)!=(k+1)k!$.
\end{merksatz}
